\documentclass[11pt, a4paper]{article}

\usepackage[utf8]{inputenc}
\usepackage[T1]{fontenc}
\usepackage[italian]{babel}
\usepackage{geometry}
\usepackage{hyperref}
\usepackage{amsmath}
\usepackage{amssymb}
\usepackage{enumitem}

\geometry{a4paper, margin=2.5cm}

\title{Gestione dei Dataset Sbilanciati nel Machine Learning}
\author{}
\date{}

\begin{document}

\maketitle

\section{Introduzione e Definizione del Problema}

Il \emph{Machine Learning} (ML) si basa sull'addestramento di algoritmi per
riconoscere pattern e prendere decisioni basate sui dati. Tuttavia, uno dei
problemi più critici e comuni nella classificazione, sia binaria sia multiclasse,
è lo \textbf{sbilanciamento dei dati}, o \emph{data imbalance}.

Un dataset è considerato sbilanciato quando la distribuzione delle classi target
è fortemente asimmetrica: una classe, detta \textbf{classe maggioritaria},
domina numericamente su un'altra, detta \textbf{classe minoritaria}. Questo
squilibrio può verificarsi anche nei problemi di regressione, quando sono
presenti valori anomali molto più alti o più bassi rispetto alla media.

\subsection{Scenari tipici di applicazione}

Esempi comuni di dataset sbilanciati si trovano nei seguenti contesti:

\begin{itemize}
    \item rilevamento di frodi bancarie, ad esempio 99\% transazioni legittime
    e 1\% frodi;
    \item diagnosi di malattie rare, ad esempio 90\% pazienti sani e 10\%
    malati;
    \item previsione dell'abbandono dei clienti, o \emph{customer churn};
    \item rilevamento di difetti di fabbricazione;
    \item cybersecurity.
\end{itemize}

\subsection{Il Paradosso dell'Accuratezza}

Gli algoritmi di machine learning sono spesso progettati per ottimizzare
l'accuratezza globale e minimizzare l'errore totale. Di conseguenza, in presenza
di dataset sbilanciati, tendono a favorire la classe maggioritaria, ignorando
quasi del tutto quella minoritaria.

Immaginiamo un sistema anti-frode dove il 99\% delle transazioni è legittimo,
cioè appartiene alla classe 0, mentre l'1\% è fraudolento, cioè appartiene alla
classe 1. Un modello mal addestrato potrebbe semplicemente decidere di prevedere
sempre la classe 0 per ogni transazione.

In questo caso, la sua accuratezza sarebbe pari al 99\%, un valore apparentemente
molto alto. Tuttavia, il modello sarebbe completamente inutile, perché non
individuerebbe mai l'1\% di frodi, che rappresenta il vero obiettivo del sistema.

L'errore sulla classe minoritaria, ad esempio non rilevare una frode o un tumore,
ha spesso un costo reale, finanziario o umano, molto più elevato rispetto a un
falso allarme sulla classe maggioritaria.

\section{Pre-elaborazione e Metriche di Valutazione}

\subsection{La Regola d'Oro: Prevenzione del Data Leakage}

Un errore critico nella gestione dei dati sbilanciati è applicare tecniche di
bilanciamento sull'intero dataset prima di dividerlo in training set e test set.

La regola fondamentale è la seguente: qualsiasi tecnica di ricampionamento deve
essere applicata esclusivamente sul training set e solo dopo aver effettuato il
\emph{train-test split}.

Inoltre, per la valutazione finale del modello, è importante utilizzare dati di
test originali e non campionati, cioè \emph{unsampled data}. In caso contrario,
il modello potrebbe soffrire di overfitting sulla distribuzione alterata e i
risultati ottenuti non rifletterebbero correttamente la realtà.

\subsection{Metriche Specifiche per Dati Sbilanciati}

In presenza di dataset sbilanciati, l'accuratezza deve essere sostituita o
affiancata da metriche basate sulla matrice di confusione. Questa permette di
analizzare:

\begin{itemize}
    \item veri positivi;
    \item falsi positivi;
    \item veri negativi;
    \item falsi negativi.
\end{itemize}

Le metriche più utili sono le seguenti:

\begin{itemize}
    \item \textbf{Precision}: indica la proporzione di previsioni positive
    corrette rispetto a tutte le previsioni positive effettuate dal modello.
    Un'alta precisione indica pochi falsi allarmi.

    \item \textbf{Recall}, o \textbf{Sensitivity}: indica la proporzione di veri
    positivi individuati rispetto a tutti i casi positivi realmente presenti nel
    dataset. È cruciale quando il costo di un falso negativo è alto.

    \item \textbf{F1-score}: è la media armonica tra precision e recall. È una
    metrica molto utile per misurare le prestazioni sulla classe minoritaria in
    contesti sbilanciati.

    \item \textbf{Specificità}: indica la proporzione di veri negativi
    individuati correttamente.

    \item \textbf{Balanced Accuracy}: corrisponde alla media tra sensibilità e
    specificità.

    \item \textbf{Curve ROC e Precision-Recall}: nei dataset sbilanciati, la
    curva Precision-Recall, o AUC-PR, è spesso più informativa della classica
    AUC-ROC. La curva ROC tratta le classi in modo più simmetrico e può risultare
    meno sensibile ai miglioramenti specifici sulla classe minoritaria.
\end{itemize}

\section{Soluzioni a Livello di Dati}

Le tecniche a livello di dati, dette anche \emph{resampling techniques}, alterano
fisicamente il dataset per rendere più equilibrata la distribuzione delle classi.

\subsection{Sottocampionamento}

Il sottocampionamento, o \emph{undersampling}, riduce il numero di istanze della
classe maggioritaria. Può essere utile per dataset di grandi dimensioni, ma
comporta il rischio di perdere informazioni preziose, fenomeno noto come
\emph{information loss}.

\subsubsection{Random Undersampler}

Il \emph{Random Undersampler} scarta campioni della classe maggioritaria in modo
casuale. Spesso riduce eccessivamente i dati, impedendo al modello di apprendere
correttamente.

\subsubsection{Tomek Links}

I \emph{Tomek Links} rappresentano un metodo selettivo. L'idea è individuare
coppie di dati appartenenti a classi opposte che sono molto vicine tra loro,
cioè situate vicino al confine decisionale.

Una volta individuata una coppia di questo tipo, viene rimossa l'istanza della
classe maggioritaria, con l'obiettivo di pulire il confine decisionale.

\subsubsection{Edited Nearest Neighbors}

Edited Nearest Neighbors, abbreviato in ENN, è una tecnica che pulisce il dataset
rimuovendo i campioni della classe maggioritaria che si trovano circondati da
vicini appartenenti alla classe minoritaria.

\subsection{Sovracampionamento}

Il sovracampionamento, o \emph{oversampling}, aumenta il numero di campioni della
classe minoritaria per avvicinarlo o eguagliarlo a quello della classe
maggioritaria.

\subsubsection{Random Oversampler}

Il \emph{Random Oversampler} duplica campioni già esistenti della classe
minoritaria. Il rischio principale è l'overfitting, poiché il modello può
memorizzare campioni ridondanti invece di apprendere regole generali.

\subsubsection{SMOTE}

SMOTE, acronimo di \emph{Synthetic Minority Over-sampling Technique}, non si
limita a duplicare i dati, ma crea nuovi dati sintetici.

L'algoritmo individua un punto della classe minoritaria e cerca i suoi $k$ vicini
più prossimi, secondo la logica dei \emph{K-Nearest Neighbors}. Successivamente,
traccia segmenti immaginari tra il punto e i suoi vicini e genera nuovi punti
sintetici posizionandoli casualmente lungo questi segmenti.

Questo procedimento aggiunge varianza e riduce il rischio di overfitting rispetto
alla semplice duplicazione dei campioni.

\subsubsection{ADASYN}

ADASYN è una tecnica simile a SMOTE, ma si concentra maggiormente sulle zone in
cui il modello ha più difficoltà a separare le classi, cioè vicino ai confini
decisionali.

\subsection{Metodi Ibridi}

I metodi ibridi, o \emph{combined resampling}, combinano oversampling e
undersampling. In questo modo si cerca di aumentare la rappresentazione della
classe minoritaria e, allo stesso tempo, pulire la classe maggioritaria.

Esempi di metodi ibridi sono:

\begin{itemize}
    \item SMOTEENN, cioè SMOTE combinato con ENN;
    \item SMOTETomek, cioè SMOTE combinato con Tomek Links.
\end{itemize}

\section{Soluzioni a Livello Algoritmico e Sensibili al Costo}

Le soluzioni a livello algoritmico non alterano direttamente il dataset. Al
contrario, modificano il modo in cui l'algoritmo apprende dai dati.

\subsection{Cost-Sensitive Learning e Pesi di Classe}

Nel \emph{cost-sensitive learning}, gli algoritmi assegnano pesi o penalità
diverse agli errori commessi sulle diverse classi.

Per esempio, invece di considerare ogni errore allo stesso modo, si può
modificare la funzione di perdita, o \emph{loss function}. Se il modello sbaglia
a prevedere una transazione normale, appartenente alla classe maggioritaria, la
penalità può essere pari a 1. Se invece sbaglia a prevedere una frode,
appartenente alla classe minoritaria, la penalità può essere pari a 10, oppure
proporzionale al grado di sbilanciamento.

In questo modo, l'algoritmo è costretto matematicamente a prestare maggiore
attenzione alla classe rara.

\subsection{Regolazione della Soglia}

Nella classificazione binaria, la soglia predefinita per assegnare un campione
alla classe positiva è spesso pari a 0.5. Se lo sbilanciamento è forte, si può
abbassare questa soglia, ad esempio a 0.3, forzando il modello a prevedere la
classe minoritaria più spesso.

Questo tende ad aumentare la recall, ma può anche aumentare il numero di falsi
positivi.

\subsection{Scelta dell'Algoritmo}

Alcuni algoritmi sono naturalmente più resistenti allo sbilanciamento. I modelli
basati su alberi, come Decision Tree e Random Forest, e i metodi ensemble, come
XGBoost e AdaBoost, possono gestire meglio lo sbilanciamento rispetto a modelli
lineari, poiché isolano le classi in rami specifici.

\section{Ensemble Methods Applicati allo Sbilanciamento}

Una tecnica avanzata consiste nella creazione di un'architettura ibrida che
combina partizionamento dei dati, o \emph{batching}, e apprendimento d'insieme,
o \emph{ensemble learning}.

Supponiamo di avere 3000 dati normali, appartenenti alla classe maggioritaria, e
1000 dati frodi, appartenenti alla classe minoritaria. L'undersampling classico
scarterebbe 2000 dati normali. L'approccio ensemble, invece, divide i 3000 dati
normali in 3 blocchi da 1000.

Si creano quindi tre modelli distinti:

\begin{itemize}
    \item Modello 1: addestrato su Blocco Normale 1 e 1000 frodi;
    \item Modello 2: addestrato su Blocco Normale 2 e 1000 frodi;
    \item Modello 3: addestrato su Blocco Normale 3 e 1000 frodi.
\end{itemize}

Infine, le previsioni dei tre modelli vengono aggregate tramite voto a
maggioranza, o \emph{majority vote}. In questo modo, l'algoritmo apprende su set
bilanciati, ma nessun dato viene completamente sprecato.

Esistono implementazioni già disponibili in Python, nella libreria
\texttt{imblearn}, come:

\begin{itemize}
    \item \texttt{BalancedRandomForestClassifier};
    \item \texttt{EasyEnsembleClassifier}.
\end{itemize}

\section{Anomaly Detection per Squilibri Estremi}

Quando la classe minoritaria è estremamente rara, ad esempio pari allo 0.001\%,
le normali tecniche di classificazione possono fallire. In questi casi, il
paradigma cambia: non si cerca più di distinguere direttamente due classi, ma si
istruisce un modello non supervisionato a definire la normalità, cioè la classe
maggioritaria.

Qualsiasi osservazione che esca da questi confini matematici viene etichettata
come anomalia, o \emph{outlier}.

Gli algoritmi più usati in questo contesto sono:

\begin{itemize}
    \item \textbf{One-Class SVM}: crea un'ipersfera attorno ai dati normali; i
    punti che cadono fuori vengono considerati anomali.

    \item \textbf{Isolation Forest}: isola le anomalie osservando che esse
    richiedono meno suddivisioni in un albero decisionale per essere separate dal
    resto dei dati.

    \item \textbf{Local Outlier Factor}, o LOF, e \textbf{Autoencoder}.
\end{itemize}

\section{Risultati Empirici e Contraddizioni tra le Fonti}

La teoria non trova sempre riscontro nella pratica. Questo mostra che nel machine
learning non esiste una regola universale valida per ogni dataset e ogni modello.

\subsection{Il Caso Controverso del Metodo SMOTE}

Dal punto di vista teorico, SMOTE viene spesso descritto come un'evoluzione utile
per prevenire l'overfitting. In alcune prove empiriche su reti neurali per il
\emph{customer churn}, SMOTE è risultato molto efficace, innalzando
drasticamente l'F1-score della classe minoritaria.

In altri esperimenti, invece, applicato a una regressione logistica per il
rilevamento frodi, SMOTE ha prodotto risultati scadenti, generando molti falsi
positivi e un F1-score molto basso.

Questo mostra che l'efficacia di SMOTE dipende fortemente dal dataset, dal
modello scelto e dal contesto applicativo.

\subsection{Pesi di Classe e Focal Loss}

L'uso automatico di pesi bilanciati, ad esempio tramite
\texttt{class\_weight = ``balanced''}, non garantisce sempre buoni risultati.
In alcuni casi può generare un numero eccessivo di falsi positivi.

Al contrario, l'impostazione manuale di pesi asimmetrici può risultare molto più
efficace, ad esempio assegnando peso 1 alla classe maggioritaria e peso 10 alla
classe minoritaria.

Anche tecniche più sofisticate, come la \emph{Focal Loss}, non garantiscono
necessariamente un miglioramento. In alcuni esperimenti, la penalizzazione
avanzata degli errori ha addirittura peggiorato l'F1-score.

\subsection{La Disillusione dei Metodi Ensemble}

L'unione tra undersampling ed ensemble learning è teoricamente interessante, ma
non sempre porta ai risultati migliori. In alcuni casi, dividere il dataset in
batch, addestrare più modelli e aggregare le previsioni con voto a maggioranza ha
prodotto un miglioramento solo modesto, inferiore rispetto a tecniche più semplici
come oversampling o SMOTE.

\subsection{Undersampling: Distruttivo o Efficace?}

Il sottocampionamento casuale è spesso criticato perché può causare una forte
perdita di informazioni. Tuttavia, metodi di undersampling più intelligenti,
come Edited Nearest Neighbors, possono risultare molto efficaci.

ENN, infatti, non elimina dati in modo casuale, ma pulisce selettivamente il
dataset, rimuovendo osservazioni problematiche vicine ai confini decisionali.
Per questo motivo, in alcuni esperimenti ha ottenuto risultati eccellenti,
raggiungendo un F1-score elevato e riducendo drasticamente i falsi positivi.

\section{Conclusione}

La gestione dei dataset sbilanciati richiede attenzione sia nella fase di
pre-elaborazione sia nella fase di valutazione del modello. L'accuratezza non è
una metrica sufficiente, perché può mascherare completamente il fallimento sulla
classe minoritaria.

Le principali strategie disponibili sono il resampling, i metodi sensibili al
costo, la regolazione della soglia decisionale, gli ensemble methods e, nei casi
più estremi, l'anomaly detection.

Tuttavia, nessuna tecnica è universalmente migliore delle altre. La scelta deve
essere guidata dal dataset, dal modello, dal costo degli errori e dalle metriche
di valutazione più adatte al problema specifico.

\end{document}